\FloatBarrier\section{Conclusão}


O desenvolvimento deste trabalho permitiu analisar, na prática, o funcionamento de uma aplicação distribuída gerenciada pelo Kubernetes, com ênfase na disponibilidade dos serviços, na recuperação automática de falhas e na escalabilidade. Os experimentos demonstraram que, diante da remoção de uma réplica da aplicação, o orquestrador foi capaz de identificar a alteração e iniciar sua substituição, restabelecendo a quantidade de três instâncias configuradas. Esse comportamento evidenciou o mecanismo de autorrecuperação (\textit{self-healing}) e a importância da manutenção do estado desejado para a continuidade do atendimento.

Os testes de escalabilidade também demonstraram a flexibilidade proporcionada pela orquestração. A expansão manual de três para dez réplicas comprovou a possibilidade de aumentar a quantidade de instâncias da aplicação por meio de uma alteração na configuração desejada. Complementarmente, o escalonamento automático, realizado com o Horizontal Pod Autoscaler (HPA), evidenciou a capacidade do Kubernetes de ajustar a quantidade de réplicas em resposta à utilização de CPU. Esses resultados permitiram compreender a diferença entre o gerenciamento manual de recursos e a adaptação automática da infraestrutura às condições de carga observadas.

A integração com o PostgreSQL e os experimentos de persistência e isolamento de rede ampliaram a compreensão sobre as dependências existentes entre os componentes de um sistema distribuído. Durante a recriação do banco de dados, a aplicação continuou respondendo às requisições HTTP, embora as operações de gravação estivessem temporariamente indisponíveis. Após a recuperação, a contagem de visitas foi retomada a partir dos valores anteriormente registrados, demonstrando a importância do armazenamento persistente. Os testes de política de rede também evidenciaram como o controle da comunicação entre os serviços contribui para a organização e a segurança da infraestrutura.

Dessa forma, a realização dos experimentos proporcionou à equipe um aprendizado sobre a gestão de sistemas interconectados. Foi possível compreender que a disponibilidade de uma aplicação não depende apenas da existência de contêineres em execução, mas também da comunicação entre os serviços, do funcionamento das dependências e da configuração adequada dos mecanismos de recuperação. Além disso, o acompanhamento dos estados dos pods, das respostas da aplicação e dos registros experimentais demonstrou a importância da observação conjunta dos componentes para identificar falhas e avaliar o comportamento do sistema.

Por fim, os objetivos propostos para a prática foram atendidos nos experimentos documentados, especialmente quanto ao acesso à aplicação pelo Service, à autorrecuperação das réplicas e à escalabilidade de três para dez instâncias. Entretanto, como os testes foram realizados em um ambiente Minikube de nó único, os resultados não permitem afirmar que a arquitetura suportaria a perda completa de uma máquina ou do armazenamento. Como possibilidade de evolução, destaca-se a realização de experimentos em clusters com múltiplos nós e a ampliação das medições de desempenho e disponibilidade. Assim, o trabalho demonstrou tanto os benefícios da orquestração com Kubernetes quanto a importância de reconhecer os limites da infraestrutura implementada.
